% Final §5 -- Conclusion and future work (plan §7.1: 0.4 pages).
\section{Conclusion and future work}\label{sec:conclusion}
SeatRelay replaces one form per session with one registration, one queue and a timed hold per
workshop, and relays every released seat automatically. Each workshop is serialised through one
FIFO message group, and every decision commits in one conditional DynamoDB transaction. On the
deployed system \SeriesPassed{} of \SeriesExperiments{} system experiments met their targets
(failure recovery without e-mail, then off on \code{dev}), no \code{verify} call found an
invariant violation, and the load test met every target up to \PerfDevMaxVUsMet{} virtual users,
also with all of them on one workshop. A first series failed for two reasons: the account's Lambda quota
of 10 (CloudWatch tied every HTTP 503 to a throttled invocation) and one late expiry-race probe; with the quota raised and the probe's send time calibrated, the same deployed code passed.
Per-offer timers with tick-first semantics take the
sweep interval off the backfill path (\IfPending{\BackfillExpiredPNinetyFiveS}{\TimerPathRelaysN{} timer-path relays on
\code{dev} so far, at most \TimerPathRelayMaxMs\,ms}{expiry-driven backfill p95 \BackfillExpiredPNinetyFiveS\,s on
\code{dev}, against 120\,s}),
and the 202 protocol is settled by bounded polling. In our usage model the cloud build costs
\usd\CostTierOneKCloudListUsd{} a month at list price for 1{,}000 registrations and is cheaper than a
self-hosted machine that pays for its line and operations up to \CostCrossoverListFullRegs{}
registrations a month. Of the hypotheses (\S\ref{sec:hypotheses}), H2 is supported in simulation and H3 partly (automatic relay
shown). The simulated research package illustrates H1 and the organiser side of H3 with
\SimInterviewN{} scripted interviews and \SimStudyN{} participants; its results are assumptions-driven,
not human evidence. The single scripted study misses the completion and SUS reference lines;
state-interpretation failures motivate clearer confirmation, ranking and date-entry guidance.

That package is complete; interviews and a study with people, empirical price validation and a
natural pilot remain open. Future work: (i)~batch commands of one
workshop per worker invocation, with a shared snapshot, for throughput; (ii)~server push only if
the user study finds the 10\,s status refresh too slow; (iii)~billing and limits for the
organiser plan; (iv)~a second region with a replicated table and identity store, if one region's
availability is not enough; (v)~the Cognito Lite tier, as Cognito dominates the cost at scale; (vi)~a mail domain of our own
with DKIM, for notifications and sign-up codes.
